\documentclass[10pt,a4paper]{amsart}
\usepackage[margin=19mm]{geometry}
\usepackage{amsmath,amssymb,mathtools}
\usepackage[scr=rsfs,cal=euler]{mathalfa}
\usepackage{xfrac}
\usepackage[unicode,hidelinks,bookmarks=false]{hyperref}
\newcommand{\ie}{i.e.\ }
\newcommand{\cf}{cf.\ }
\newcommand{\resp}{respectively\ }
\newcommand{\Subsection}{\subsection}
\providecommand{\nobreakdash}{\nobreak\hbox{-}}
\setlength{\emergencystretch}{2em}
\multlinegap0pt
\usepackage{amssymb}
\usepackage{mathtools}
\usepackage[shortlabels]{enumitem}
\usepackage{xcolor}
\newtheorem{proposition}{Proposition}
\newcommand{\Gal}{\operatorname{Gal}}
\title{Counterexamples to O'Neil's Period-Index Problem on Elliptic Curves}
\author{Xiaoguang Shang, Cheng Niu}
\date{Source: arXiv:2608.29288v1 (CC BY 4.0)}
\begin{document}
\maketitle

\noindent\emph{Source and licence.} Counterexamples to O'Neil's Period-Index Problem on Elliptic Curves. Version arXiv:2608.29288v1 (CC BY 4.0), distributed by arXiv under the Creative Commons Attribution 4.0 International licence, http://creativecommons.org/licenses/by/4.0/, as recorded in the arXiv licence metadata for this exact version and shown on the abstract page https://arxiv.org/abs/2608.29288v1. Full licence notice: \texttt{licenses/arxiv-2608.29288-CC-BY-4.0.txt}.\par\smallskip

\noindent\textit{Editorial scope.} Counted unit: the article's proposition lem:strong-commutator (source0.tex lines 612--620). The article's complete original proof of it is delivered with it (source0.tex lines 622--657, one \texttt{proof} environment of 36 lines). No mathematics is written, repaired or re-derived: the statement and the proof are the article's own lines, in the article's own notation, and no retained line is substituted or re-flowed.  No further source-local context is needed: every cross-reference of every retained line resolves inside the two delivered spans.  Nothing was omitted from any retained span, so the package carries no omission record.  No retained line contains a citation, so the wrapper carries no bibliography and no bibliography entry.  No retained line contains a figure environment, an image-inclusion command or drawing code, so no image is delivered and the package declares no asset dependency.  Editorial formatting only, in addition: an AMS \texttt{amsart} preamble that restates the article's own declarations and macros used by the retained lines, namely the environments \texttt{proposition} on separate counters, together with the standard packages those lines need.  The retained environments print in this document as Proposition 1 in order of appearance, whereas the article prints the same items with the numbering of its own full document.  The complete native source of the article as distributed in the arXiv e-print for this exact version is bundled unchanged as \texttt{sources/208838/source0.tex}.
\begin{proposition}
	\label{lem:strong-commutator}
	There exists a central commutator
	\[
	c\in\Gal(A/K)
	\]
	whose restriction to $F_{\mathrm{MW}}$ is the identity and whose
	restriction to $H$ is equal to $2P$.
\end{proposition}

\begin{proof}
	Let	$J=H\cap F_{\mathrm{MW}}$. Since $\Gal(F_{\mathrm{MW}}/L)$ is an elementary abelian
	$2$-group, the element $2Q\in\Gal(H/L)$ acts trivially on $J$.
	Therefore, the compatible pair	$(2Q,\operatorname{id})$
	defines an element $ s\in\Gal(A/L)$.
	
	Let $\widetilde{\tau}\in\Gal(A/K)$	be any lift of $\tau\in\Gal(L/K)$, and define
$
	c=\widetilde{\tau}s\widetilde{\tau}^{-1}s^{-1}.
$
	It is immediate that $c|_{F_{\mathrm{MW}}}=1$.
	On the other hand, restricting to $H$, we obtain
	\[
	c|_H
	=
	(\widetilde{\tau}s\widetilde{\tau}^{-1})|_H(s^{-1})|_H
	=
	(2P+2Q)-2Q
	=
	2P.
	\]
Finally, let $g\in \Gal(A/K)$. On $F_{\mathrm{MW}}$, we have
\[
(gcg^{-1})|_{F_{\mathrm{MW}}}
=
\operatorname{id}
=
c|_{F_{\mathrm{MW}}}.
\]
On the other hand, restricting to $H$, since $2P$ is $K$-rational, we have
\[
(gcg^{-1})|_H=c|_H=2P.
\]
Therefore, $gcg^{-1}=c$, and hence $c$ is central in
$\Gal(A/K)$. This completes the proof.
\end{proof}

\end{document}
